\section{Limitations}

The current study is intentionally scoped to geometry-checkable relation
families. Relations whose correctness depends primarily on functional, social,
affordance, or language-pragmatic context are outside the H001 metric claim.
The evaluation also uses closed-set/GT-object relation rows and exact
predicate-label recall, so the results should be read as relation-reliability
evidence rather than full open-vocabulary scene graph generation performance.

The Open3DSG experiment is a reproduced averaged-BLIP variant with explicit
filtering and coverage caveats. The checkpoint is selected by train-dev
validation loss before H001 held-out inspection, the training and validation
splits are filtered to preprocessed-ready subgraphs, and H001 evaluation covers
the loadable Open3DSG scope. These choices do not invalidate the scoped
reliability result, but they prevent claims of exact non-averaged Open3DSG
reproduction or broad open-vocabulary SOTA.

The visual and qualitative evidence should also be read with the correct scope.
The reduced visual sanity check and Open3DSG qualitative case inspection are
reviewer-defense and failure-mechanism evidence, not large-scale independent
human audits. They help identify failure modes and residual calibration risk,
but the quantitative claim remains tied to the fixed metric denominator and
Docker-generated tables.

Finally, \pgeom{} is a calibrated reliability score, not a proof of physical
correctness. Residual high-confidence but rule-violated cases motivate future
work on broader predicate families, stronger calibration, larger visual audits,
modern VLM relation sources, and downstream tasks where relation reliability can
be tested through planning, navigation, alignment, or embodied reasoning.

